%==================================================================
% SAMPUL DALAM  (Bahasa Inggris — mengikuti contoh hal. 2)
% Variabel Bahasa Inggris didefinisikan di a0-identitas.tex
%==================================================================
\thispagestyle{empty}
\begin{center}
    \begin{spacing}{1.5}
        \textbf{\large \MakeUppercase{\judulen}}\\[1.5cm]
    \end{spacing}

    \textbf{\MakeUppercase{\labeltype}}\\[0.3cm]

    \ifdefempty{\submitlineen}{}{\textbf{\submitlineen}\\[1.5cm]}

    % Logo diatur dengan HEIGHT (bukan width) agar tinggi sampul tetap
    % sama untuk logo berbentuk bujur sangkar maupun memanjang.
    \includegraphics[height=3.2cm]{logo-institusi}\\[1.5cm]

    \textbf{Submitted by:}\\
    \textbf{\MakeUppercase{\penulis}}\\
    \textbf{\nim}

    \vfill

    \textbf{\MakeUppercase{\prodien}}\\
    \textbf{\MakeUppercase{\fakultasen}}\\
    \textbf{\MakeUppercase{\universitasen}}\\
    \textbf{\MakeUppercase{JAKARTA}}\\
    \textbf{\tahun}
\end{center}
\clearpage
\setcounter{page}{2}
